% Independently audited in full by root and B after the four precision repairs.
\subsection{Eight excess arrows with acyclic support}

\begin{lemma}[The rank-one form of the small-corner lemma]
\label{C:small-corner-two-stalks-rank-one}
Let a finite category $S$ be the source block of a collage
whose sink corner $N$ has at most four elements.  Suppose
every cross stalk is nonempty and has at most two elements.
If the collage has rational second-cohomology dimension at
most one, its degree-two square vanishes.  The opposite
statement holds for a small source corner.
\end{lemma}
\begin{proof}
This is the homological scope of the proof of
Lemma~\ref{C:small-corner-two-states-per-object}.
For a projective trivial right $\mathbb QN$-module, the
actual orbit quotient is a rational ring equivalence to
an objectwise one- or two-state categorical act cone.
All its positive products vanish by
Lemma~\ref{C:two-states-objectwise-act-cups}.
Otherwise the minimal ideal of $N$ is a right-zero band
of size two or three.  Its retractions coincide on each
one- or two-element stalk.  Their common image is an actual
biset retract on which $N$ acts trivially.  Since $BN$ is
contractible, that core collage is equivalent to an
objectwise one- or two-state categorical act cone and
has zero products.  The relative quotient is a suspension.
If the core has nonzero second cohomology, its retract
pullback spans the assumed rank-one space; otherwise every
degree-two class lifts from the suspension quotient.
In both cases the square vanishes.  These arguments use
only the stated rank bound, not bounded homology or finite
CW type of the collage.
\end{proof}

\begin{lemma}[Two parallel-arrow excesses with small corners]
\label{C:two-parallel-small-corners}
Let $C$ be a connected finite category with acyclic support,
all corner orders at most four, support cycle rank at most
six, and total parallel-arrow excess at most two.
If $\dim H^2(BC;\mathbb Q)\leq1$, its square vanishes.
\end{lemma}
\begin{proof}
All its corners are rationally acyclic by the small-monoid
classification.  If there is at most one parallel hom-set,
or if its two parallel hom-sets cannot occur serially,
Lemma~\ref{C:parallel-hom-cut} applies.
Thus there are exactly two doubled hom-sets
$C(a,b)$ and $C(c,d)$ with $b\leq c$.

If $b<c$, let $A$ be the full downset of $c$.
Every hom-set from $A\setminus\{c\}$ to $c$ is a singleton:
neither of the two doubled hom-sets ends at $c$.
Thus $BA$ is the join of the predecessor subcategory with
the corner at $c$, and has zero second rational cohomology.
The complementary full subcategory has only singleton
cross hom-sets and has zero second cohomology by
Lemmas~\ref{C:singleton-support-H2} and
\ref{C:flag-second-cycle-budget}, including its isolated
corners.  The two-block suspension lemma proves the claim.

It remains that the two doubled hom-sets are
$C(a,b)$ and $C(b,c)$.  Suppose $b$ has a neighbor
$x>b$ other than $c$.  Let $A$ be the full downset of $b$.
All hom-sets $A\to x$ are singleton, so the full
subcategory on $A\cup\{x\}$ is a join and has zero
second cohomology.  Every global degree-two class
therefore restricts to zero on $A$.  The complement
of $A$ has only singleton cross hom-sets and zero
second cohomology as above.  The actual two-block
quotient again kills the square.  A neighbor $x<b$
other than $a$ gives the opposite argument with the
upset of $b$.

Otherwise $b$ has precisely the two neighbors $a,c$.
Let $T$ be the full triangle on $a,b,c$ and let
$D=C\setminus\{b\}$.  Their actual nerve union is all
of $BC$, and their intersection is the singleton edge
$a\to c$.  Its classifying space is the join of two
rationally acyclic corner spaces, hence rationally acyclic.
The subcategory $D$ has only singleton cross hom-sets and
zero second cohomology.  Mayer--Vietoris identifies
$H^2(BT;\mathbb Q)$ with the global group, and makes
fourth-cohomology restriction injective.
The triangle has hom-set sizes $(2,2,1)$.  Its last
corner has order at most four and both incoming stalks
have size at most two.  The preceding lemma kills its
rank-one square; the restriction to $D$ is already zero.
This proves the assertion.
\end{proof}

\begin{theorem}\label{C:acyclic-arrow-excess-eight}
Let $C$ have $k$ objects, $n$ arrows, and acyclic support.
If $n\leq2k+7$, then $BC\not\simeq\mathbb CP^2$.
In particular every acyclic seven-object category with
at most twenty-one arrows is excluded.
\end{theorem}
\begin{proof}
Assume $BC\simeq\mathbb CP^2$, so $C$ is connected.
The cases $k\leq2$ have at most eleven arrows and are
already excluded.  Use the notation $D,P,R$ of
Theorem~\ref{C:acyclic-arrow-excess-seven}.
That theorem treats $D+P+R\leq7$, so it remains to
consider equality
\[
 D+P+R=8.
\]
As before the support functor is onto on fundamental
groups, and its support order complex has zero first
homology over $\mathbb F_2$.

If $P=0$, the actual support functor pulls back every
degree-two class and its target has zero fourth homology
by Lemma~\ref{C:flag-second-cycle-budget}.
If the corners are semisimple and $R\leq5$, the relative
bar has strict-object length at most three and $H_4=0$.
For semisimple corners and $R\geq6$, one has
$D+P\leq2$.  The support has at most one five-object
chain.  Choose a face of that chain deleting a vertex
which is not an endpoint of any parallel hom-set used
by its top simplex, nor a nontrivial corner.
Such a vertex exists: for $D=0$ at most two parallel
hom-sets exclude at most four vertices; for $D=1$ there
is at most one parallel hom-set and one nontrivial
corner, excluding at most three vertices.

The selected face retains all varying arrow labels.
Its remaining composition involves only singleton
hom-sets, except possibly a fixed singleton product
landing in a larger hom-set.  In that event the unique
source basis vector is sent to one target basis vector,
so this factor is still an injection.  The resulting
component of the top boundary is injective.  If the one nontrivial
semisimple corner occurs, it is retained by this face;
the component is the identity on the retained balanced
tensor factors, tensored with the just-described basis
injection.  Exact semisimple coinvariants preserve this
injection.  There is no other five-object
chain to contribute to the same boundary component.
Thus the degree-four relative-bar boundary is injective
and $H_4=0$ in all semisimple cases.

We may now assume $D\geq2$ and a nonsemisimple corner.
If $P=1$, the parallel-cut and exceptional-leaf argument
of the excess-seven proof applies unchanged.  There is
at most one corner with nonzero $H^2$, since $D\leq7$.
A singleton incident edge kills its global restriction.
Without such an edge it is a leaf on the sole doubled
hom-set, whose products vanish universally.  Its other
endpoint has order at most four and is rationally acyclic,
and the remaining singleton-hom category has zero $H^2$
since $R\leq5$.  Fourth-cohomology restriction finishes
this case.  Henceforth $D,P\geq2$, so $R\leq4$.

\paragraph{Tree support.}
When $R=0$, zero second cohomology at every corner gives
the height-one suspension obstruction.  There is otherwise
a unique corner $M$ with nonzero $H^2$, of order at most
seven.  A singleton incident edge again kills all global
vertex restrictions, so assume every incident edge is parallel.

If every other corner has order at most two, group all
edges incident with $M$ into its central star.  All remaining
pieces have semisimple diagonal and one strict crossing;
they have no homology in degrees at least two, and meet
the star at rationally acyclic corners.  The star inherits
bounded rational homology and second-cohomology rank at
most one.  Exact leaf coinvariants and identification of
the trivial leaf points give an ordinary act cone over
$M$, changing only first homology.  The bounded rank-one
act lemma for $|M|\leq9$ kills its square.

The only other budget is $D=6$, $P=2$, with the positive
corner the five-element sphere monoid and one other corner
of order three.  If the positive corner has degree two,
both incident hom-sets have size two.  Their individual
collages have zero products universally, and their common
five-element corner has $H_3(-;\mathbb Q)=0$.
Restriction in degree four is therefore injective on
their union.  Every remaining piece attaches at a
rationally acyclic corner and has zero products.
If the positive corner has degree one, its hom-set has
size two or three.  Size two is universal.  Size three
is either an ordinary act over the two-row sphere monoid,
or its three-state biset with the order-three corner.
For the latter, a projective trivial module at the small
opposite corner gives an ordinary act after exact
coinvariants.  The only nonprojective order-three corner
is the identity-adjoined two-element right-zero band,
and Corollary~\ref{C:P5-right-zero-three-all-states}
gives zero products universally.  No actual
projective-plane type of the piece is assumed.
This closes tree support.

\paragraph{One independent cycle.}
Let $R=1$.  The support consists of its one triangle
and bridge attachments.  If all corners have order at
most five, their third rational homology is zero, so
fourth-cohomology restriction to these pieces is injective.
When all corners have order at most four the intersections
are rationally acyclic: each piece inherits bounded homology
and $b_2\leq1$.

For $D\leq3$, bridge products vanish by the two-block
lemma, and triangle pieces are handled by the same
endpoint zero-$H^2$ cut or middle-order-at-most-four
bounded-Tor argument as in the excess-seven proof.
For $D=4$ or $5$, suppose first all corners have order
at most four.  Here respectively $P=3$ or $2$.
A singleton adjacent hom-set gives a join block and
a zero-$H^2$ cut.  Otherwise the adjacent sizes are
at least two.  The budget leaves either three hom-set
sizes at most two, or $(2,3,1)$ or $(3,2,1)$.
In each case one extreme corner has both incident
stalks of size at most two.  Lemma~\ref{C:small-corner-two-stalks-rank-one}
kills the inherited rank-one square.  Bridge products
vanish since both their corners have zero $H^2$.

If a five-element corner occurs, all other corners
have order at most two and $P\leq3$.
A bridge at a positive five-element corner is, after
exact opposite coinvariants, an ordinary act over the
two-row sphere monoid, so has zero products.  A bridge
whose corners have zero $H^2$ uses the two-block lemma.
At the middle of the triangle, collapse the semisimple
extremes rationally.  A singleton adjacent set deletes
the middle object.  Otherwise the adjacent sizes are
at least two; either all three sets have size two,
giving an objectwise two-state cone, or the long set
is singleton and one adjacent size is two, giving
a two-row or two-column band by triangle compression.
These all have zero products.

At an endpoint, a five-element corner with zero $H^2$
is cut from the opposite semisimple two-object block.
If its $H^2$ is nonzero, it is the sphere monoid.
A singleton incident hom-set is a fixed arrow; the
opposite semisimple block has zero $H^2$, so the
fixed-arrow lemma kills all products.  Without such
a singleton, both incident sizes are at least two.
The opposite endpoint is semisimple.  Except when
the long set has size three, its two incoming sets
have size at most two after exact endpoint orbits,
so the objectwise-two-state lemma applies.
The exceptional sizes are $(2,1,3)$, up to opposite;
they force $D=4$, so the two other corners are trivial.
Delete the middle object toward the other trivial
endpoint and apply the ordinary two-row act theorem.

The only corner of order six has $D=5$, $P=2$, with
all other corners trivial.  To avoid an unjustified
corner-intersection argument at its possibly nonzero
$H_3$, group every simplex using this corner into one
piece $X$: include all its incident bridge edges and,
if it lies on the support triangle, that triangle too.
The remaining bridge pieces and, if not already included,
the unique triangle attach to this grouped piece and
to one another only at trivial corner spaces.
Hence $X$ inherits bounded rational homology and
$b_2\leq1$, and fourth-cohomology restriction is
injective.  All other pieces have trivial corners
and dimension at most two, so have zero fourth homology.

If the order-six corner is not on the triangle, $X$
is a star, and identification of its trivial leaves
gives the ordinary bounded rank-one act obstruction.
If it is in the middle of the triangle, it has no
further neighbor: an extra comparable neighbor would
give a second support cycle.  A singleton adjacent
hom-set deletes the middle, while two parallel
adjacent sets must both have size two and the long
set is singleton.  The latter triangle is an
objectwise two-state cone and has zero products.
If the order-six corner is an extreme of the triangle,
all its additional neighbors are on that same side.
When the opposite adjacent hom-set is singleton,
delete the trivial middle object toward the other
trivial extreme, leaving a star and hence an ordinary
act after identifying leaves.  Otherwise that opposite
set is parallel.  The budget $P=2$ forces one of the
two hom-sets incident with the order-six corner in
the triangle to be singleton.  This is a fixed arrow,
and its full complement in $X$ has trivial corners
and height at most one, so has zero $H^2$.
The fixed-arrow lemma finishes this case.

\paragraph{Two parallel-arrow excesses and at least two cycles.}
Suppose $P=2$ and $R\geq2$.  Then $D\leq4$.
If all corners have order at most four,
Lemma~\ref{C:two-parallel-small-corners} applies.
Otherwise $D=4$, $R=2$, and the sole nontrivial corner
is a five-element monoid $M$.  If its $H^2$ is zero,
the parallel-cut proof applies whenever the parallel
hom-sets are not adjacent; in the adjacent degree-two
case its triangle has a trivial extreme and two-point
stalks, so the same proof applies.  Alternatively one
may use the following direct positive-corner argument
only when its $H^2$ is nonzero.

If $H^2(BM;\mathbb Q)\ne0$, then $M$ is the sphere
monoid.  A singleton incident hom-set makes every
global degree-two restriction to $BM$ zero; the same
parallel-cut proof then works, with the binary triangle
exception again an objectwise two-state cone.
Without such a singleton, $M$ has one or two neighbors.
Let $X$ be the full subcategory on $M$ and these neighbors,
and let $D'$ be its complement with only the $M$ object
removed.  Their nerves cover the category and meet in
a full subcategory on at most two trivial objects.
Their third intersection homology is zero.
The category $D'$ has trivial corners and support cycle
rank at most two, so its nerve has dimension at most two.
If $M$ has one neighbor, $X$ is an ordinary two- or
three-state act over the sphere monoid and has zero
products.  If it has two incomparable neighbors, $X$
is its two-leg, two-state star; identifying the two
trivial endpoints gives an ordinary four-state act
and preserves products in degrees at least two.
If its neighbors are comparable, $X$ is a binary
triangle with a trivial extreme and hence is an
objectwise two-state cone.  Its products vanish in
every case, and restriction in degree four finishes.

\paragraph{The remaining two-cycle budget.}
For $R=2$ and $P\geq3$, the possibilities are
$(D,P)=(2,4)$ or $(3,3)$.
For $D=2$ there is just one order-three corner.
Its trivial modules have projective dimension at
most one, so every edge intersection has zero
homology in degrees at least three.  Fourth-cohomology
restriction to the triangle pieces is injective,
and they inherit bounded high homology.  At an endpoint
the order-three corner gives total bar degree at
most three; in the middle the small-middle Tor
obstruction gives zero fourth homology.  This closes
the $D=2$ case.

For $D=3$ all corners are rationally acyclic and
there is at most one corner of order at least three.
Bridges and individual triangles are handled as in
the one-cycle $D\leq3$ case.  Only a diamond remains.
Write its shared edge as $u<v$ and its exclusive
vertices as $a,b$.
When $u<a,b<v$, avoiding a zero-$H^2$ star cut forces
both potentially nonzero edges to meet at the one
larger exclusive corner, say $a$.
The shared-edge endpoints are semisimple, so that
edge has no homology in degrees at least two.
Restriction in degree four to the two triangles is
injective, and bounded high homology passes to them.
The triangle through $a$ has middle order at most
four and semisimple endpoints; the other is semisimple.
Both have zero fourth homology.

Suppose instead $a,b<u<v$, the opposite case being
dual.  If the shared edge has zero $H^2$, use the
cut $\{a,b\}\mid\{u,v\}$.  Otherwise that edge is
parallel and touches the unique larger corner.
A clean exclusive vertex, with both hom-sets to
$u,v$ singleton, gives a join and kills every global
restriction to the shared edge, proving the same cut.
Avoiding such a clean vertex uses all three parallel
excesses: the shared edge is doubled, and one of
the two incident hom-sets at each of $a,b$ is doubled.
Thus every hom-set in the diamond has size at most two.

If the larger corner is $u$, the final corner at $v$
is semisimple.  Its exact orbit collapse makes it
trivial and preserves the two-point stalk bound;
the objectwise-two-state lemma kills all products.
Suppose the larger corner is $v$.
The other three corners are semisimple, and their
full height-one source block has zero $H^2$.
The corner at $v$ has order at most four and also
has zero second rational cohomology.  The actual
two-block suspension lemma for
$\{a,b,u\}\mid\{v\}$ therefore kills every
degree-two product directly.  This excludes the
diamond and completes the two-cycle case.

\paragraph{The last three-cycle budget.}
The only case not yet covered with $R=3$ is
$D=2$, $P=3$.  There is one order-three corner.
As above, edge intersections have zero homology in
degrees at least three, and triangle pieces have zero
fourth homology after inheriting boundedness.
For a complete four-object block, all hom-set sizes
at most two give the objectwise-two-state argument
at a trivial extreme.  A single parallel hom-set
of size four gives the parallel-hom cut.
The only other pattern consists of one hom-set of
size three and another of size two.
If they are not serial, use the parallel-hom cut.
If they are separated serially, cut at the downset
of the second source, exactly as in
Lemma~\ref{C:two-parallel-small-corners}.
If they are adjacent, their common vertex has an
additional neighbor in the four-object chain;
the join-cut argument in that lemma applies.
Thus all products vanish on the four-object block,
and fourth-cohomology restriction finishes the case.

The only four-cycle budget was $D=P=2$, already
covered by the two-parallel lemma.  These alternatives
exhaust $D,P\geq2$ and $D+P+R=8$.
Every one contradicts the projective-plane square,
proving the theorem and its seven-object consequence.
\end{proof}
